\chapter{Data Understanding and Versioning}\label{ch:data}

\section{Environment}
\textbf{What and why.} I created an isolated virtual environment so the exact library versions
of this project cannot clash with other projects on my machine, and so anyone can rebuild the
same environment from \file{requirements.txt}. Every version is pinned (e.g.\ TensorFlow 2.20.0,
Keras 3.15.1, MLflow 2.22.5, Keras Tuner 1.4.7, DVC 3.67.1).

\textbf{How.} \cmd{py -3.11 -m venv .venv} creates a private copy of the interpreter and an empty
\cmd{site-packages}; activating it puts that interpreter first on \texttt{PATH}.
\cmd{pip install -r requirements.txt} installs the pinned set and \cmd{pip install -e .} installs
my own package \texttt{visionguard} (in \file{src/}) in editable mode so every script imports the
same code. \file{scripts/00_env_check.py} prints the user, host, Python, every library version,
whether the interpreter is inside a virtualenv, and what TensorFlow can see:
\cmd{tf.config.list_physical_devices("GPU")}. On my laptop TensorFlow sees only the CPU (native
Windows TensorFlow has no GPU support after 2.10), so the heavy training runs were executed on a
Colab T4 GPU with the same pinned tracking libraries (\cref{sec:2-0}).

\shot{s1_1_env}{Exact library versions, virtualenv check and GPU visibility (\file{scripts/00_env_check.py}).}

\textbf{Why MLflow 2.x and not 3.x?} MLflow 3 logs models through a new ``logged models'' REST
API. DagsHub's hosted tracking server implements the MLflow 2 API, so I pinned the client to the
latest 2.x release to keep \texttt{log\_model} and the Model Registry working against DagsHub.

\section{Dataset audit}
\textbf{What.} \file{scripts/01_audit.py} walks every file under \file{data/raw} and records, per
image: class (from the top-level folder name), file size, width, height, colour mode, format, a
SHA-256 hash of the bytes, whether the file can be fully decoded, and the error if not. The full
table is saved to \file{reports/audit/audit.csv}; summaries are written straight into this report.

\textbf{How each check works.}
\begin{itemize}
\item \emph{Readability}: \texttt{PIL.Image.verify()} catches broken headers and \texttt{load()}
forces a full pixel decode, which catches truncated files that open fine but fail later.
I deliberately use the \emph{same} decode function as training and serving
(\file{src/visionguard/imaging.py}), so ``readable'' in the audit means ``usable by the model''.
\item \emph{Duplicates}: two files are exact duplicates when their SHA-256 hashes are equal (same
bytes). I keep the first copy only. Duplicates matter because a copy in the training set and
another in the test set leaks the answer and inflates test accuracy.
\item \emph{Label conflicts}: a hash that appears under \emph{both} classes is label noise; all
copies are dropped.
\item Images smaller than 32 pixels on a side are flagged as unusable.
\end{itemize}

\gentable{audit_counts}{Files per class: total, readable, corrupt (cannot be decoded), exact duplicates and too-small images.}
\gentable{audit_stats}{Minimum, median and maximum of width, height and file size per class (readable files only).}
\gentable{audit_formats}{Image formats and colour modes per class.}

The dataset has \vgNFiles{} files: \vgNAI{} AI-generated and \vgNReal{} real. \vgNCorrupt{}
files could not be decoded and \vgNDup{} are exact duplicates (\vgNDupGroups{} duplicate groups,
\vgNConflicts{} hashes appearing in both classes). After cleaning, \vgNClean{} images remain and
are split 70/15/15 with stratification (same AI/real ratio in every split), using a fixed seed:

\gentable{splits}{Stratified train / validation / test split after removing corrupt files, duplicates and label conflicts.}
\gentable{audit_corrupt}{Files that could not be decoded (first 15) and the decoder's error.}

\shot{s1_2_audit}{Terminal output of the audit script.}

\begin{wrongbox}[Problem I had to design around]
The obvious decoder, \texttt{tf.io.decode\_image}, cannot read WEBP files, CMYK JPEGs or
palette PNGs with transparency, all of which occur in scraped data. Using it would have marked good
images as ``corrupt'' and thrown them away, so the audit, training and serving all use one Pillow
function that converts every image to RGB (compositing transparency onto white so transparent
pixels do not turn black).
\end{wrongbox}

\section{Visual EDA}
\autofig{eda_samples}{Random samples per class (top two rows AI-generated, bottom two rows real), with original size and format.}
\textbf{Interpretation.} The AI-generated images are mostly polished, centred compositions with
smooth gradients, glossy lighting and an ``illustration'' look, while the real images include
everyday clutter, uneven exposure and natural noise. Both classes contain people, animals and
landscapes, so the model cannot solve the task by topic alone; it has to use style and texture.

\autofig{eda_histograms}{Width, height, file size (log scale) and aspect ratio per class.}
\textbf{Interpretation.} The median AI image is \vgMedWAI{}$\times$\vgMedHAI{} pixels and
\vgMedKBAI{}\,KB, versus \vgMedWReal{}$\times$\vgMedHReal{} pixels and \vgMedKBReal{}\,KB for real
images. \vgSquarePctAI{}\% of AI images are exactly square against \vgSquarePctReal{}\% of real
ones, because generators output fixed square resolutions. These differences are a
\emph{shortcut risk}: a network could learn ``square and small file $\Rightarrow$ AI'' instead of
learning generation artefacts. Resizing every image to the same square input removes the size
signal from the pixels, but compression strength still differs; I return to this in the error
analysis (\cref{sec:2-6}).

\yours{After running \file{02_eda.py} on the real data, check that the two paragraphs above still
match your figures (e.g.\ which class is larger, how strong the square-image peak is) and adjust a
sentence if not.}

\section{Data version control with DVC}
\subsection*{Toy exercise}
\textbf{What and why.} In an empty folder I created a 3-row \file{scores.csv}, versioned it with
DVC, changed it, travelled back to the old version, and round-tripped it through remote storage
on DagsHub (repository \texttt{23L-0764\_A2\_Ops-dvc-toy}). Commands, in order
(\file{labs/dvc_toy.sh}):
\begin{enumerate}
\item \cmd{git init}, \cmd{dvc init}: Git tracks code and small text; DVC adds \file{.dvc/}
(config + local cache) on top of Git.
\item \cmd{dvc add scores.csv}: moves the content into the content-addressed cache
\file{.dvc/cache/<md5>} and writes the pointer \file{scores.csv.dvc}; it also adds
\file{scores.csv} to \file{.gitignore} so Git never stores the data itself. I commit the pointer.
\item Append two rows, \cmd{dvc add} again, commit: the pointer now holds a different md5.
\item \cmd{git checkout HEAD~1 -- scores.csv.dvc} restores the old pointer and
\cmd{dvc checkout} replaces the workspace file with the cached content for that md5: the 3-row
version is back. Git chooses \emph{which} version; DVC supplies the \emph{bytes}.
\item \cmd{dvc remote add -d origin s3://dvc} plus \cmd{endpointurl} points DVC at DagsHub's
S3-compatible storage; the token is set with \cmd{--local}, which writes it to
\file{.dvc/config.local} (git-ignored) instead of the committed \file{.dvc/config}.
\item \cmd{dvc push --all-commits} uploads the data for every committed version. I then deleted
\file{scores.csv} \emph{and} the local cache, and \cmd{dvc pull} downloaded it again.
\end{enumerate}

\shot{s1_4_toy_pointer}{The pointer file and the two commits (v1, v2).}
\shot{s1_4_toy_checkout}{Returning to the older version with \cmd{git checkout} + \cmd{dvc checkout}.}
\shot{s1_4_toy_pushpull}{Push to DagsHub, delete local data and cache, pull it back.}
\shot{s1_4_toy_dagshub}{The toy repository on DagsHub with the DVC-tracked file.}

\subsection*{What the pointer file holds, and why data does not belong in Git}
The \file{.dvc} file is a few lines of YAML: the \texttt{md5} of the content (for a directory, the
md5 of a manifest listing every file and its hash), the total \texttt{size}, the number of files
and the \texttt{path}. It is a fingerprint, not data. Git is built for small text that changes
line by line: every version of every file is stored forever in \file{.git}, so committing
hundreds of megabytes of images makes every clone download all of it, binary diffs are
meaningless, and GitHub rejects files over 100\,MB. With DVC, Git history stays tiny, and the
data version is the md5 in the pointer, which I also log to MLflow as the tag
\texttt{dvc\_data\_version}, so every experiment says exactly which data it used.

\subsection*{Applying it to VisionGuard}
\cmd{dvc add data/raw} tracks the whole dataset directory as one unit (\file{data/raw.dvc}),
with DagsHub storage of the \texttt{23L-0764\_A2\_Ops} repository as the default remote. The
split lists (\file{data/splits/*.csv}, a few KB) are committed to Git directly, because they
are small text and must change together with the code that produced them.

\shot{s1_4_dataset_dvc}{\cmd{dvc add data/raw}, commit of the pointer, and \cmd{dvc push} of the dataset.}
\shot{s1_4_dagshub_data}{DagsHub repository page showing \file{data/raw} tracked by DVC.}

\begin{wrongbox}[Problems to watch for]
A plain \cmd{dvc push} uploads only the data of the \emph{current} commit, so after deleting the
cache an older version could not be restored from the remote; \cmd{dvc push --all-commits}
uploads every version referenced in Git history. The second trap is credentials:
\cmd{dvc remote modify} without \cmd{--local} writes the token into \file{.dvc/config}, which is
committed, so the token would be published.
\end{wrongbox}
